\section{Обучение устойчивой модели вознаграждения (RRM): предотвращение взлома функции вознаграждения}

\subsection{Ключевые идеи}

Основная цель статьи~\cite{liu2025rrmrobustrewardmodel}~-- устранить зависимость модели вознаграждения от артефактов, не связанных с контекстом запроса, и научить модель оценивать исключительно смысловую релевантность ответа.

Авторы рассматривают каузальную модель зависимости предпочтений человека $C$ от контекстуального сигнала качества $S$ и артефакта $A$.  
Традиционные модели вознагражения обучаются по парам $(x, y_w, y_l)$ и не могут отделить влияние $S$ и $A$, что приводит к ложной зависимости $A \rightarrow C$.

Для устранения этого эффекта вводится \emph{причинно-следственная структура}, обеспечивающая независимость:
$$
A \perp C \mid (Y_1, Y_2), \qquad A \perp C \mid S
$$

Для реализации этой независимости используется метод аугментации данных на основе перестановок инструкций и ответов из разных примеров.  
Сбалансировав артефакты между выбранными и отклонёнными ответами, модель утрачивает возможность угадывать предпочтения на основе формата и становится устойчивой к взлому.

\subsection{Математическое описание метода}

Рассмотрим множество пар $(x, y_1, y_2)$, где $x$~-- запрос, а $y_1$ и $y_2$~-- два ответа, один из которых предпочтён человеком.  
Цель обучения модели вознаграждения~-- аппроксимировать вероятность того, что ответ $y_1$ предпочтительнее ответа $y_2$ для данного запроса $x$.  
Эта вероятность моделируется через сигмоидальную функцию:

\begin{equation}
P(y_1 \succ y_2 \mid x) = \sigma(r(x, y_1) - r(x, y_2)),
\end{equation}

где $r(x, y)$~-- функция вознаграждения, определяющая качество ответа $y$ в контексте $x$,  
а $\sigma(z) = \frac{1}{1 + e^{-z}}$~-- сигмоидальная функция активации.  

\emph{Традиционная функция потерь} для модели вознаграждения имеет вид:

\begin{equation}
\mathcal{L}_{\text{обычн.}} = - \mathbb{E}_{(x, y_w, y_l)} 
\left[ \log \sigma(r_\phi(x, y_w) - r_\phi(x, y_l)) \right],
\end{equation}

где $y_w$~-- выбранный (предпочтённый) ответ, $y_l$~-- отклонённый,  
а $r_\phi$~-- функция вознаграждения с параметрами $\phi$.  
Такое обучение учитывает только реальные пары предпочтений и не отделяет артефакты, не зависящие от контекста запроса.

Для устранения этой зависимости вводится \emph{аугментированное множество данных}, включающее перестановки ответов из разных примеров:

\begin{equation}
\mathcal{D}_{\text{аугм.}} = 
\Bigl\{ \bigl(x^{(i)}, y^{(\sigma_1(i))}_w, y^{(\sigma_2(i))}_l \bigr) 
\Bigr\},
\end{equation}

где $\sigma_1, \sigma_2$~-- перестановки индексов, выбираемые из других обучающих примеров,  
что обеспечивает балансировку артефактов между выбранными и отклонёнными ответами.  
Таким образом, модель не может полагаться на формат или длину ответа как на сигнал предпочтения.

Для \emph{нейтральных пар}, где оба ответа не связаны с запросом или имеют одинаковое качество, вводится «ничейный» сигнал предпочтения:

\begin{equation}
y_w \approx y_l \quad \Rightarrow \quad 
\sigma(r(x, y_w) - r(x, y_l)) \approx 0.5,
\end{equation}

что эквивалентно требованию нулевого градиента для различия между ответами без контекстуального преимущества.

Итоговая функция потерь устойчивой модели вознаграждения объединяет исходное и аугментированное обучение:

\begin{equation}
\mathcal{L}_{\text{уст.}} = 
\mathcal{L}_{\text{обычн.}} 
+ \lambda \, 
\mathbb{E}_{(x, y_w, y_l) \in \mathcal{D}_{\text{аугм.}}}
\left[ 
\log \sigma\bigl(r_\phi(x, y_w) - r_\phi(x, y_l)\bigr)
\right],
\end{equation}

где $\lambda$~-- коэффициент балансировки между исходными и аугментированными данными.  

Таким образом, оптимизация минимизирует общую функцию потерь:
\begin{equation}
\min_{\phi} \; \mathcal{L}_{\text{уст.}},
\end{equation}
что приводит к параметрам $\phi^\ast$, обеспечивающим устойчивость к артефактам, не зависящим от контекста запроса.

\subsection{Алгоритм}

\begin{enumerate}
  \item \emph{Сбор данных.}  
  Собираются пары предпочтений $(x, y_w, y_l)$, где $x$~-- запрос, $y_w$~-- выбранный ответ, $y_l$~-- отклонённый ответ.  

  \item \emph{Аугментация данных.}  
  На основе исходных пар создаются дополнительные пары для балансировки артефактов:
  \begin{itemize}
    \item \emph{Неконтекстуальные пары}~-- запрос $x$ используется совместно с ответами $(y^{(\sigma_1(i))}_w, y^{(\sigma_2(i))}_l)$ из других примеров, что разрушает связь между артефактами и контекстом;
    \item \emph{Нейтральные пары}~-- ответы не связаны с запросом или имеют одинаковое качество, что обеспечивает нулевой градиент предпочтения.
  \end{itemize}

  \item \emph{Обучение модели вознаграждения.}  
  Модель вознаграждения $r_\phi(x, y)$ обучается на объединённом множестве исходных и аугментированных данных.  
  Минимизируется итоговая функция потерь:
  \begin{equation}
  \mathcal{L}_{\text{уст.}} = 
  \mathcal{L}_{\text{обычн.}} 
  + \lambda \, 
  \mathbb{E}_{(x, y_w, y_l) \in \mathcal{D}_{\text{аугм.}}}
  \left[ 
  \log \sigma(r_\phi(x, y_w) - r_\phi(x, y_l))
  \right],
  \end{equation}
  где $\lambda$~-- коэффициент балансировки.

  \item \emph{Обучение политики.}  
  Устойчивая модель вознаграждения используется для обучения политики генерации $\pi_\theta(y|x)$ методом прямой оптимизации предпочтений (Direct Preference Optimization, DPO~\cite{rafailov2024direct}).  
  Функция потерь DPO имеет вид:
  \begin{equation}
  \mathcal{L}_{\text{DPO}} =
  -\mathbb{E} \left[
  \log \sigma \left(
  \beta 
  \log \frac{\pi_\theta(y_w|x)}{\pi_{\text{ref}}(y_w|x)}
  - 
  \beta 
  \log \frac{\pi_\theta(y_l|x)}{\pi_{\text{ref}}(y_l|x)}
  \right)
  \right],
  \end{equation}
  где $\pi_{\text{ref}}$~-- эталонная (базовая) политика, $\beta$~-- коэффициент масштабирования, контролирующий степень отклонения от эталонной модели.

  \item \emph{Оценка качества.}  
  Полученная политика проверяется по метрикам RewardBench, MT-Bench и AlpacaEval-2 для оценки точности предпочтений и согласованности длины ответов.
\end{enumerate}

\subsection{Наборы данных}

Для обучения и проверки устойчивой модели вознаграждения использовались наборы данных из коллекции \emph{RLHFlow}~\cite{li2024rlhflow}.  
Эта коллекция объединяет несколько открытых источников предпочтений человека, широко применяемых при обучении с подкреплением по обратной связи от человека (RLHF).

\begin{itemize}
  \item \emph{HH-RLHF} (\emph{Helpful and Harmless Reinforcement Learning from Human Feedback})~\cite{bai2022training}~-- набор диалоговых данных, в котором аннотаторы оценивают, насколько ответы моделей полезны и безопасны.  
  Используется для формирования предпочтений между более и менее корректными ответами на пользовательские запросы.  

  \item \emph{SHP} (\emph{Stanford Human Preferences})~\cite{ethayarajh2022shp}~-- корпус сравнений пар коротких текстов из форумов Reddit, где пользователи оценивают, какой ответ является более подходящим.  
  Этот набор хорошо подходит для обучения моделей предпочтений в неформальных диалоговых контекстах.  
  

  \item \emph{HelpSteer}~\cite{wang2023helpsteer}~-- набор инструкций и ответов, аннотированных по степени полезности (от 0 до 4).  
  Применяется для обучения модели вознагражения оценивать качество ответов по шкале, а не только по бинарным предпочтениям.  
  

  \item \emph{UltraFeedback}~\cite{cui2023ultrafeedback}~-- один из крупнейших наборов данных предпочтений, включающий сравнения ответов, сгенерированных большими языковыми моделями.  
  Используется для настройки моделей на высокоуровневые характеристики качества: глубину рассуждения, креативность, связность и корректность.  
\end{itemize}

После применения процедур аугментации объём данных увеличен примерно в 14 раз.  
Итоговый обучающий набор содержит около $2.4 \times 10^6$ пар сравнений, обеспечивающих устойчивое распределение артефактов по всем подмножествам.

\subsection{Метрики}

Для количественной и качественной оценки работы устойчивой модели вознаграждения и обученных политик использовались три метрики, предложенные в работах по оценке моделей согласования с человеческими предпочтениями.

\begin{itemize}
  \item \emph{RewardBench accuracy}~\cite{liu2024rewardbench}~-- доля корректных парных сравнений, выполняемых моделью вознаграждения.  
  RewardBench представляет собой стандартизированный бенчмарк.  
  Метрика вычисляется как доля случаев, когда модель правильно определяет предпочтённый ответ из пары $(y_w, y_l)$, согласно человеческим аннотациям.

  \item \emph{MT-Bench score}~\cite{zheng2023judging}~-- экспертная (человеческая или модельная) оценка качества ответов по десятибалльной шкале.  
  Для оценки используется многоуровневый тест из 80 инструкций, охватывающих области рассуждения, кода, математики и диалога.  
  Метрика применяется для сравнения генеративных политик, согласованных с человеческими предпочтениями.

  \item \emph{AlpacaEval-2 length-controlled win rate}~\cite{dubois2024length}~-- процент побед модели при сравнении пар ответов одинаковой длины.  
  В отличие от исходного AlpacaEval, данная версия контролирует смещение по длине, устраняя эффект, при котором более длинные ответы необоснованно получают более высокие оценки.  
  Метрика измеряет устойчивость модели к артефактам длины и отражает реальное качество сгенерированных ответов.
\end{itemize}

\subsection{Модели}

Для обучения и экспериментов использовались современные модели и методы оптимизации, предназначенные для согласования языковых систем с человеческими предпочтениями.  

\begin{itemize}
  \item \emph{Базовая модель вознаграждения.}  
  В качестве основы использована языковая модель \emph{Gemma-2-9b-it}~\cite{team2024gemma}, разработанная компанией Google DeepMind.  
  Эта модель содержит около $9$ миллиардов параметров и прошла предварительное обучение на корпусе инструктивных диалогов с последующим дообучением (инструкционным тонким тюнингом).  
  Используется как инициализация для построения функции вознаграждения $r_\phi(x, y)$, оценивающей предпочтительность ответа $y$ при запросе $x$. 

  \item \emph{Оптимизатор.}  
  Для настройки параметров модели применён оптимизатор \emph{AdamW}~\cite{loshchilov2019adamw}, обеспечивающий стабильное обучение при больших размерах батча и эффективное регуляризующее действие за счёт корректного весового распада.

  \item \emph{Метод обучения политик.}  
  Для обучения политики генерации ответов $\pi_\theta(y|x)$ на основе устойчивой модели вознаграждения применён метод \emph{прямой оптимизации предпочтений} (Direct Preference Optimization, DPO)~\cite{rafailov2024direct}.  
  Этот метод минимизирует функцию потерь, напрямую выражающую предпочтения человека, без использования явного обучения с подкреплением.  
  Такой подход обеспечивает большую устойчивость и воспроизводимость результатов по сравнению с традиционным RLHF.

  \item \emph{Методы сравнения.}  
  Для оценки качества RRM выполнено сравнение с методами \emph{ODIN}~\cite{chen2023odin} и длино-контролируемой прямой оптимизацией предпочтений~\cite{park2024lengthdpo}.  
  Метод ODIN реализует двухканальную архитектуру, где один выход оценивает качество, а другой контролирует влияние длины ответа.  
  Длино-контролируемая оптимизация (Length-controlled DPO) вводит штрафной член в функцию потерь, снижающий смещение в пользу длинных ответов.  
  Эти методы служат сильными базовыми линиями для проверки устойчивости модели к взлому функции вознаграждения.
\end{itemize}

\subsection{Результаты}

\begin{itemize}
  \item Точность на RewardBench увеличена с $80.61\%$ до $84.15\%$;
  \item Средний балл MT-Bench вырос с $7.27$ до $8.31$;
  \item Контролируемый по длине win-rate на AlpacaEval-2 повысился с $33.46\%$ до $52.49\%$.
\end{itemize}
